% PDF p. 221
\BackSection{习题}
\begin{enumerate}[itemsep=.5\baselineskip]
  \item \textbf{习题 4.1}\par\nopagebreak
  假设
  \begin{itemize}
    \item $M=\R^2$，坐标为 $(x,y)$。
    \item $N=\R^2$，坐标为 $(u,v)$。
    \item 映射 $f\colon\R^2\to\R^2$ 定义为
    \[
      u=y^2,\qquad v=y+x^2.
    \]
    \item $N$ 上的 1-形式 $\omega=u\,dv$。
  \end{itemize}
  计算 $f^*\omega$。

  \item \textbf{习题 4.2}\par\nopagebreak
  设
  \begin{itemize}
    \item \textbf{1-形式} $\omega=x\,dx+y\,dy$（径向场）。
    \item \textbf{向量场} $X=-y\,\partial/\partial x+x\,\partial/\partial y$（$\R^2$ 中的旋转向量场）。
  \end{itemize}
  计算 $\mathcal L_X\omega$。

  \item \textbf{习题 4.3}\par\nopagebreak
  在 $\R^3$ 中，设坐标为 $(x^1,x^2,x^3)$，且
  \begin{itemize}
    \item $X=x^3\,\partial/\partial x^1-x^2\,\partial/\partial x^2$（旋转向量场）。
    \item $Y=\partial/\partial x^3$（沿 $x^3$ 方向的常向量场）。
  \end{itemize}
  计算 $(\mathcal L_X Y)^i$。

  \item \textbf{习题 4.4}\par\nopagebreak
  在 $\R^2$ 中，设坐标为 $(x_1,x_2)$，且
  \begin{itemize}
    \item $X=x^2\,\partial/\partial x^1+x^1\,\partial/\partial x^2$。
    \item $\omega=x^1\,dx^1-x^2\,dx^2$。
  \end{itemize}
  计算 $(\mathcal L_X\omega)_i$。

  \item \textbf{习题 4.5}\par\nopagebreak
  给出计算 $(2,1)$ 型张量场李导数 $(\mathcal L_X T)^{i_1,i_2}_{j}$ 的一般公式。

  \item \textbf{习题 4.6}\par\nopagebreak
  请证明，在球面 $S^2(R)$ 上，
  \[
    \Delta_g f=\frac1{R^2}\left(f_{\theta\theta}+\cot\theta\,f_\theta
                   +\frac1{\sin^2\theta}f_{\varnothing\varnothing}\right).
  \]

  \item \textbf{习题 4.7}\par\nopagebreak
  请证明，在双曲平面 $(x+iy,y>0)$ 上，
  \[
    \Delta_g f=y^2(f_{xx}+f_{yy}).
  \]
\end{enumerate}

% PDF p. 222
\BackSection{参考文献}
\begin{bookreferences}
  \item Eliasof, M., Haber, E. and Treister, E. (2024). Graph neural reaction giffusion models. \href{https://arxiv.org/abs/2406.10871}{arXiv:2406.10871}.
  \item Gilmer, J., Schoenholz, S. S., Riley, P. F., Vinyals, O. and Dahl, O. E. (2017). Neural message passing for quantum chemistry. \href{https://arxiv.org/abs/1704.01212}{arXiv:1704.01212}.
  \item Gogoshin, G. and Rodin, A. S. (2023). Graph neural networks in cancer and oncology research: Emerging and future trends. \emph{Cancers} 15(24): 5858.
  \item Jiang, D., Wu, Z., Hsieh, C.-Y., et al. (2021). Could graph neural networks learn better molecular representation for drug discovery? A comparison study of descriptor-based and graph-based models. \emph{Journal of Cheminformatics} 13: 12.
  \item Kengkanna, A. and Ohue, M. (2024). Enhancing property and activity prediction and interpretation using multiple molecular graph representations with MMGX. \emph{Communications Chemistry} 7: 74.
  \item Kobayashi, S. and Nomizu, K. (1969). \emph{Foundations of Differential Geometry} (Volume I \& II). John Wiley \& Sons, New York.
  \item Lee, J. M. (2012). \emph{Introduction to Smooth Manifolds} (Volume 218). Springer, New York.
  \item Martín-Landrove, M. (2017). Reaction-diffusion models for Glioma tumor growth. \href{https://arxiv.org/abs/1707.09409}{arXiv:1707.09409}.
  \item Oates, C. J., Dondelinger, F., Bayani, N., Korola, J., Gray, J. W. and Mukherjee, S. (2014). Causal network inference using biochemical kinetics. \href{https://arxiv.org/abs/1406.0063}{arXiv:1406.0063}.
  \item Proietti, M., Ragno, A., La Rosa, B., Ragno, F. R. and Capobianco, R. (2024). Explainable AI in drug discovery: self-interpretable graph neural network for molecular property prediction using concept whitening. \emph{Machine Learning} 113: 2013–2044.
  \item Spivak, M. (1995). \emph{Calculus on Manifolds}. Addison-Wesley Publishing Company, New York.
  \item Spivak, M. (2008). \emph{A Comprehensive Introduction to Differential Geometry} (Volumes 1–5). Publish or Perish, Inc.
  \item Tu, L. W. (2011). \emph{An Introduction to Manifolds}. Second Edition. Textbook. Springer, New York.
  \item Zhao, B., Xu, W., Guan, J. and Zhou, S. (2024). Molecular property prediction based on graph structure learning. \emph{Bioinformatics} 40: btae304. \url{https://doi.org/10.1093/bioinformatics/btae304}
\end{bookreferences}
